\SourceLeaf{374}%
von \(AO\) ein Stück \(OM=NL\) ab und zieht \(ML\); wenn man ferner in \(M\) an \(MA\) abermals einen Winkel von der Grösse des Winkels \(A\) setzt, dessen anderer Schenkel \(MC\) die \(AN\) in dem Punkte \(C\) zwischen \(A\) und \(L\) trifft, hierauf von \(AM\) ein Stück \(MB=LC\) abschneidet und \(BC\) zieht, und sodann diese Construction auf ähnliche Art fortsetzt, so dass auf der Linie \(OA\) die Punkte \(O,M,B,E,G,K\) u.\,s.\,w., auf der Linie \(NA\) hingegen die Punkte \(N,L,C,D,F,H\) u.\,s.\,w. liegen, so wird behauptet, dass die Stücke \(OM,MB,BE,EG,GK\) u.\,s.\,w. oder die ihnen resp. gleichen \(NL,LC,CD,DF,FH\) u.\,s.\,w. eine abweichende Progression bilden.

Den Beweis dieses Lehrsatzes sucht der Verf. apagogisch so zu führen, dass er die übrigen möglichen Fälle, wenn der Lehrsatz nicht wahr wäre, aufzählt, und die Unstatthaftigkeit eines jeden zu erweisen versucht. Der Verf. behauptet nemlich, dass unter jener Voraussetzung einer von folgenden fünf Fällen Statt finden müsste. Die auf einander folgenden Stücke, von \(OM\) an gerechnet, wären
\par\noindent\hspace*{\parindent}1) alle einander gleich, oder
\par\noindent\hspace*{\parindent}2) jedes nachfolgende grösser als das vorhergehende, oder
\par\noindent\hspace*{\parindent}3) einige einander gleich und das darauf folgende grösser oder kleiner, oder
\par\noindent\hspace*{\parindent}4) einige auf einander folgende nähmen fortschreitend ab, und die darauf folgenden fortschreitend zu oder
\par\noindent\hspace*{\parindent}5) sie würden abwechselnd grösser und kleiner.
\par\noindent In dieser Aufzählung ist der mögliche Fall übergangen, dass die Stücke anfangs fortschreitend zu und dann fortschreitend abnähmen, und nach Rec. eigener Überzeugung (deren tiefer liegende Gründe hier aber nicht angeführt werden können) wäre dessen Erledigung gerade die Hauptsache und die eigentliche Auflösung des Gordischen Knotens. Inzwischen kann man zugeben, dass diese Auslassung hier in so fern wenig auf sich hat, als die Beweisart des Verf. für die Unstatthaftigkeit des dritten Falles, wenn sie zulässig wäre, auch auf diesen Fall von selbst erstreckt werden könnte. Allein eben diesem angeblichen Beweise der Unstatthaftigkeit des dritten Falls können wir keine Gültigkeit zugestehen. Der Verf. stellt die Sache so vor. Wenn z.\,B. in dem dritten Falle angenommen wird, die beiden ersten Stücke seien gleich, das dritte aber grösser, so wäre \(DC\) also grösser als \(CL\). Da nun aber \(AML\) gleichfalls ein gleichschenkliges Dreieck ist, dem dieselbe Grundbedingung zukommt, wie dem ursprünglichen Dreieck \(AON\), so müsste, wenn jener dritte Fall mit seiner angenommenen Unterabthei%
